\section{Proof Details}
\label{app:proofs}

\subsection{Conditional validity of the trial pseudo-outcome}

Fix \(q=(q_0,q_1)\) and condition on \(X\).  Under trial randomization and
consistency,
\begin{align}
\mathbb E_R\left[\frac{A}{e(X)}\{Y-q_1(X)\}\mid X\right]
&=\mathbb E_R\{Y(1)-q_1(X)\mid X\},
\\
\mathbb E_R\left[\frac{1-A}{1-e(X)}\{Y-q_0(X)\}\mid X\right]
&=\mathbb E_R\{Y(0)-q_0(X)\mid X\}.
\end{align}
Substituting these two equalities into the definition of \(\varphi\) cancels the
two fitted regressions and yields
\(\mathbb E_R\{\varphi(V;q)\mid X\}=\tau(X)\).  Applying the identity to
\(q=\widehat\mu_R\), \(q=\widehat\mu_O\), and their affine blend gives
\begin{align}
\mathbb E_R(Z_\lambda\mid X)&=\tau(X),
&
\mathbb E_R(\Delta\mid X)&=0.
\end{align}
These statements are conditional on the pre-evaluation sigma-field so that the
fitted functions are fixed.

\subsection{Proof of Theorem~\ref{thm:ate-variance}}

Using \(\mathbb E_R(Z_\lambda)=\theta\) and exact transport,
\begin{align}
\mathbb E_O\{r(X)g(X)\}=\mathbb E_R\{g(X)\},
\end{align}
the conditional mean of Equation~\eqref{eq:ate-estimator} is \(\theta\).  Rewrite
the centered estimator as
\begin{align}
\widehat\theta_r(\lambda,\omega)-\theta
&=
(\mathbb P_{R,n}-P_R)\{Z_\lambda-\omega g\}
+(\mathbb P_{O,N}-P_O)\{\omega rg\}.
\end{align}
The two summands are independent conditional on \(\mathcal F\).  The variance of
an empirical mean of \(m\) independent observations is the population variance
divided by \(m\), which proves Equation~\eqref{eq:ate-variance}.  Unbiased sample
variances on the two independent evaluation blocks therefore give the displayed
fixed-selected-coefficient variance estimator in Section~\ref{sec:ate}.  Coverage
of its Wald interval additionally uses the stated conditional Lindeberg and
moment assumptions; variance unbiasedness alone is not a central limit theorem.

If an estimated ratio \(\widehat r\) replaces \(r\), the conditional mean changes
by exactly
\begin{align}
\omega P_O\{(\widehat r-r)g\}.
\end{align}
Thus honest evaluation does not restore exact unbiasedness.  Negligibility of
both the root-\(n\) drift and the OBS empirical remainder is required before the
fixed-ratio evaluation variance can be used as a first-order variance formula.

\subsection{Proof of Theorem~\ref{thm:joint-oracle}}

Expand the first variance in Equation~\eqref{eq:ate-variance} using
\(Z_\lambda=Z_0+\lambda\Delta\).  After collecting terms and multiplying by
\(n\), the coefficient-dependent part is
\begin{align}
2C\lambda-2D\omega+A\lambda^2-2E\lambda\omega+B\omega^2,
\end{align}
which is \(2\ell^\top\beta+\beta^\top\mathbf H\beta\).  The representation
\begin{align}
\begin{pmatrix}u&v\end{pmatrix}\mathbf H\begin{pmatrix}u\\v\end{pmatrix}
&=\operatorname{Var}_R(u\Delta-vg)
+\frac nN v^2\operatorname{Var}_O(rg)
\end{align}
proves positive semidefiniteness.  If \(AB-E^2>0\), the first-order condition
\(\mathbf H\beta=-\ell\) has the unique solution
\begin{align}
\lambda^\star&=\frac{ED-BC}{AB-E^2},
&
\omega^\star&=\frac{AD-EC}{AB-E^2}.
\end{align}
When \(\mathbf H\) is singular, nonnegativity of the variance quadratic implies
\(\ell\in\operatorname{Range}(\mathbf H)\); otherwise a null-space direction
would make the expression arbitrarily negative.  The complete minimizer set is
therefore \(-\mathbf H^+\ell+\operatorname{Null}(\mathbf H)\).  Continuity of the
quadratic proves attainment on every prespecified compact feasible set.

Finally, because \(g\) is \(X\)-measurable,
\begin{align}
E
&=\mathbb E_R[\{g(X)-P_Rg\}\mathbb E_R(\Delta\mid X)]=0,
\\
D
&=\operatorname{Cov}_R\{\mathbb E_R(Z_0\mid X),g(X)\}
=\operatorname{Cov}_R\{\tau(X),g(X)\}.
\end{align}
Substitution yields Equation~\eqref{eq:axis-oracle}.  Since the population
objective is minimized over a set containing \((0,0)\), its oracle value cannot
exceed the trial-only value.  This comparison concerns the population objective,
not coefficients estimated from a finite tuning sample.

\subsection{Proof of Corollary~\ref{cor:joint-oracle-variance}}

Write the conditional variance as
\begin{align*}
V(\beta)
&=
V_0+
\frac1n
\{2\ell^\top\beta+\beta^\top\mathbf H\beta\},
&
V_0
&=
\frac1n\operatorname{Var}_R(Z_0).
\end{align*}
When \(AB-E^2>0\), the matrix \(\mathbf H\) is positive definite and the unique minimizer is \(\beta^\star=-\mathbf H^{-1}\ell\). Substitution gives
\begin{align*}
V(\beta^\star)
&=
V_0-
\frac1n\ell^\top\mathbf H^{-1}\ell.
\end{align*}
For \(\mathbf H=\left(\begin{smallmatrix}A&-E\\-E&B\end{smallmatrix}\right)\),
\begin{align*}
\mathbf H^{-1}
&=
\frac1{AB-E^2}
\begin{pmatrix}B&E\\E&A\end{pmatrix},
\\
\ell^\top\mathbf H^{-1}\ell
&=
\frac{BC^2-2ECD+AD^2}{AB-E^2}.
\end{align*}
Positive definiteness makes this quadratic form nonnegative and makes it zero exactly when \(\ell=0\). Setting \(E=0\) yields \(C^2/A+D^2/B\).

If \(\mathbf H\) is singular, boundedness below of the variance objective implies \(\ell\in\operatorname{Range}(\mathbf H)\). Evaluating at the minimum-norm solution \(-\mathbf H^+\ell\) gives
\begin{align*}
V_{\min}
&=
V_0-
\frac1n\ell^\top\mathbf H^+\ell.
\end{align*}
The pseudoinverse is positive definite on \(\operatorname{Range}(\mathbf H)\), so the reduction is strictly positive exactly when \(\ell\neq0\). Finally, continuity gives attainment on every nonempty compact feasible set \(\mathcal C\). If \(0\in\mathcal C\), the constrained minimum is at most \(V(0)=V_0\), and it is strictly smaller exactly when some feasible \(\beta\) makes \(2\ell^\top\beta+\beta^\top\mathbf H\beta<0\).

\subsection{Prediction-powered CATE risk identity}

Fix \(t\) and condition on \(X\).  Since
\(\mathbb E_R(Z_\lambda-\tau(X)\mid X)=0\),
\begin{align}
\mathbb E_R\{Z_\lambda-t(X)\}^2
&=
\mathbb E_R\{Z_\lambda-\tau(X)\}^2
+\mathbb E_R\{t(X)-\tau(X)\}^2.
\end{align}
Exact transport also gives
\begin{align}
\mathbb E_O\{r(X)(g(X)-t(X))^2\}
&=\mathbb E_R\{(g(X)-t(X))^2\}.
\end{align}
The correction in Equation~\eqref{eq:dr-fusion-loss} therefore has mean zero,
which proves Equation~\eqref{eq:cate-target}.  In particular, population expected
loss is flat in \(\omega\).  The part depending on \(\lambda\) is
\begin{align}
Q(\lambda)
&=Q(0)+2C\lambda+A\lambda^2,
\end{align}
so when \(A>0\) its minimizer is \(-C/A\), the same integrated-noise oracle as
in the population-orthogonal ATE problem.  If \(A=0\), then \(\Delta=0\) almost
surely by conditional mean zero, and \(Q\) is flat in \(\lambda\).

For a fixed \(t\) and \(\lambda\), the empirical loss is the sum of independent
empirical means of \(U_{t,\lambda}-\omega G_t\) under the trial and
\(\omega rG_t\) under the observational law.  Expanding its variance as a
quadratic in \(\omega\) and differentiating gives
Equation~\eqref{eq:loss-variance-omega}.  Because \(t\) is fixed in this
calculation, the formula does not identify the coefficient minimizing the risk
of a learner chosen after seeing the same loss.

For the linear sieve, expand Equation~\eqref{eq:dr-fusion-loss} in \(\beta\), add
\(\rho\beta^\top P\beta\), and set the gradient to zero.  The resulting normal
equation is
\begin{align}
\{(1-\omega)H_R+\omega H_{O,r}+\rho P\}\beta
&=h_\lambda+\omega d_g,
\end{align}
which proves Equation~\eqref{eq:linear-sieve} whenever the matrix is invertible.

\subsection{Loss-targeted balance and approximate drift}

Suppose a nonnegative out-of-sample-evaluable weight \(w\) satisfies
\begin{align}
P_O\{w(g-t)^2\}=P_R\{(g-t)^2\}
\quad\text{for every }t\in\mathcal T.
\end{align}
Replacing \(r\) by \(w\) then leaves the population correction exactly zero for
every candidate, so the target identity remains valid.  If
\(t_\beta=b^\top\beta\), expanding
\begin{align}
(g-b^\top\beta)^2
&=g^2-2\beta^\top bg+\beta^\top bb^\top\beta
\end{align}
shows that balancing \(bg\) and \(bb^\top\) preserves the minimizer.  Balancing
\(g^2\) is needed only to preserve the absolute loss value.

For approximate weights, the difference between the weighted and target
population losses is bounded uniformly over \(\mathcal T\) by
\(|\omega|\Delta_\mathcal T(w)\).  Let \(t_w\) minimize the weighted population
loss and \(t_0\) minimize the target loss.  Applying the uniform bound at both
minimizers gives
\begin{align}
L_0(t_w)
&\leq L_w(t_w)+|\omega|\Delta_\mathcal T(w)
\\
&\leq L_w(t_0)+|\omega|\Delta_\mathcal T(w)
\\
&\leq L_0(t_0)+2|\omega|\Delta_\mathcal T(w),
\end{align}
which is the claimed excess-risk bound.  If the same finite sample is used both
to impose exact sieve balance and to fit \(\beta\), the balanced linear and
quadratic coefficient terms cancel from the empirical correction (and the
constant cancels as well if \(g^2\) is balanced).  This algebra explains why
weight learning and candidate fitting require honest separation or a separate
outer-fold argument.
